Cyber-physical critical infrastructure relies on networked sensing, estimation, and control. This dependence creates an attack surface through which corrupted measurements can alter both the estimated state and the resulting control action. False-data injection attacks (FDIAs) have been studied extensively in power, water, and other infrastructure systems \cite{liu2011false,liang20162015,clark2017protecting,brentan2021cyber}. A particularly important class of FDIAs exploits knowledge of the estimator or detector to produce a significant physical effect while maintaining a small detection statistic \cite{mo2010false}.

Model-based bad-data detection commonly evaluates an innovation or residual against a prescribed threshold. Examples include generalized likelihood-ratio tests, Bayesian detectors, and chi-square tests associated with Kalman filtering \cite{tang2016detection,ding2018survey,mo2010false}. These methods offer interpretable statistical thresholds, but their performance depends on the plant and noise models. Data-driven classifiers provide a complementary mechanism for identifying compromised channels and have been applied to FDIA detection and localization \cite{ozay2016machine,boyaci2021joint,esmalifalak2011stealth}. Their effectiveness, however, depends strongly on whether the training data represent attacks near the decision boundary.

Adversarial data generation addresses this limitation by enriching the training distribution. Existing approaches include optimization-based attack synthesis and generative models trained from attack corpora \cite{yang2013false,mohammadpourfard2020generation,ahmadian2018cyber}. The former can be computationally demanding, whereas the latter presuppose a representative collection of feasible attacks. A physics-guided generator provides an alternative: known plant and estimator dynamics can be used as fixed discriminators, allowing attack samples to be generated without first collecting a labeled attack dataset.

Attack localization alone does not ensure resilient estimation. A localization model may retain compromised measurements or discard measurements that are important for observability. Secure and resilient estimators therefore exploit attack sparsity, sensor redundancy, and support information \cite{fawzi2014secure,pajic2014robustness,pajic2015attack,zheng2021Resilient}. Prior work has investigated $\ell_1$ estimators, support-weighted estimators, and resilient Kalman filters \cite{lee2015secure,zheng2020attack}. A remaining design question is how localization reliability should be translated into a sensor subset whose cleanliness, observability, and estimation error can be quantified simultaneously.

\emph{Contributions}: This paper addresses that question in a Kalman-filter setting. First, finite-horizon maps are derived from a sparse measurement attack to the closed-loop critical output and to the normalized Kalman-filter innovation. These maps define physics-based discriminators for generating attacks subject to an explicit chi-square detection-probability constraint. Second, a channel-wise localization network is combined with a calibrated sensor-selection rule. Under stated conditional-calibration and independence assumptions, the rule guarantees that the particular subset supplied to the estimator is attack-free with a prescribed probability. Third, the analysis relates the realized attack sparsity, the observability threshold, the selection reliability, and the conditional subset-KF covariance. This relation quantifies the estimation penalty incurred when higher selection reliability requires fewer measurements.

The remainder of the paper is organized as follows. Section~\ref{sec:notation} introduces the notation. Section~\ref{sec:system_model} presents the stochastic CPS model, Kalman filter, and chi-square detector. Section~\ref{sec:attack_generation} derives the attack-response maps and formulates physics-guided attack generation. Section~\ref{sec:detection_localization} develops channel-wise localization and calibrated sensor selection. Section~\ref{sec:resilient_estimator} analyzes the resulting subset Kalman filter. Section~\ref{sec:simulation} presents the water-distribution case study, and Section~\ref{sec:conclusion} concludes the paper.
